\section{Introduction}
\label{sec:introduction}

Modern data-management systems increasingly modify query results
in order to satisfy external requirements such as fairness,
representation, diversity, consistency, and data-quality constraints.
Examples include fairness-aware query refinement systems that modify
query predicates to satisfy demographic constraints
\cite{ERICA,MK09},
provenance-based frameworks for bias analysis and mitigation
\cite{LMSJ23},
and consistent query answering (CQA) systems that repair
inconsistent databases before evaluating queries
\cite{arenas1999,bertossi2006,bertossi2011}.
Although these approaches successfully enforce the desired constraints,
they often introduce substantial semantic changes whose consequences
are difficult for users to understand.

To illustrate the problem, consider the following aggregate query over
the Adult dataset. The query computes the average capital gain among
highly educated individuals who work long hours. \ag{Revise queries to have AVG or SUM, add example tuples for the dataset, highlight the changes in the query, and place in an example enviroment.} \alon{Addressed: income- based aggregation, example tuples and highlighting the changes in the query}
\begin{example}[Query refinement over the Adult dataset]
Consider the following simplified records from the Adult dataset:

\begin{center}
\resizebox{\columnwidth}{!}{%
\begin{tabular}{c|c|c|c|c|c}
id & gender & education & occupation & hoursPerWeek & income \\
\hline
1 & Male   & Bachelors & Exec-managerial & 60 & 1 \\
2 & Male   & Bachelors & Prof-specialty  & 55 & 1 \\
3 & Female & Masters   & Prof-specialty  & 40 & 0 \\
4 & Female & Masters   & Adm-clerical    & 35 & 0 \\
5 & Male   & Masters   & Exec-managerial & 52 & 1 \\
\end{tabular}%
}
\captionof{table}{Simplified Adult records. The binary attribute
\texttt{income} equals 1 for individuals earning more than 50K and 0
otherwise.}
\end{center}

The original query computes the fraction of high-income individuals
among highly educated employees who work long hours:

\begin{lstlisting}
SELECT AVG(income) AS highIncomeRate
FROM Adult
WHERE education IN ('Bachelors','Masters')
  AND hoursPerWeek >= 50;
\end{lstlisting}

Suppose that the population selected by this query violates a
representation constraint requiring at least 25\% female records.
A query-refinement mechanism may therefore produce the following
refined query, where the modified predicate values are shown in bold:

\begin{lstlisting}[escapeinside={(*@}{@*)}]
SELECT AVG(income) AS highIncomeRate
FROM Adult
WHERE education = (*@\textbf{'Masters'}@*)
  AND hoursPerWeek >= (*@\textbf{30}@*);
\end{lstlisting}

The refined query changes the population over which the aggregate is
computed. In the simplified records above, ids 1 and 2 may be removed
because their education level is Bachelors, while ids 3 and 4 may be
added because the hours-per-week threshold was relaxed. Thus, the
refinement does not merely change the aggregate value; it changes the
underlying population whose income distribution is being summarized.

For example, the refinement may increase female representation from
19.5\% to 30.6\%, thereby satisfying the representation constraint.
However, understanding why the refinement succeeds is more difficult
than verifying that the constraint is satisfied. Did the refinement
simply introduce more women into the result? Did it alter the
occupational composition of the selected population? Did it change the
income distribution? Did it affect particular demographic groups?
\end{example}

As we show later, the refinement succeeds not merely because it
increases female representation, but because it replaces a
population dominated by high-income workers with intensive work
schedules by a substantially different population with distinct
occupational and economic characteristics.
Understanding and communicating such semantic changes is the
central goal of this work.

Although existing refinement systems provide the modified query,
they offer limited support for explaining the semantic consequences
of the refinement \cite{ERICA,MK09}.

The same challenge arises beyond query refinement.
In consistent query answering, the answer produced from a repaired
database may differ substantially from the answer obtained from the
original inconsistent database
\cite{arenas1999,bertossi2006,bertossi2011}.
Similarly, data-cleaning pipelines, provenance-based bias-mitigation
frameworks, and fairness-aware analytics frequently generate multiple
related query outputs whose semantic differences must be understood
and justified \cite{LMSJ23,MLJ22}.

A natural solution is to provide explanations describing the
semantic differences between two query outputs.
However, generating such explanations is challenging.

First, multiple semantic factors may simultaneously contribute
to the observed difference between two query outputs.
Changes in demographic composition, education levels, occupations,
income distributions, or family characteristics may all contribute
to the observed transformation.
Identifying the most informative factors is therefore non-trivial.

Second, explanations must characterize the \emph{difference}
between two outputs rather than the properties of either output
in isolation.
Many characteristics may be common to both outputs and therefore
irrelevant to understanding the observed change.

Third, explanations are often derived from sensitive demographic,
financial, medical, or employment data.
Releasing exact explanation statistics may inadvertently reveal
information about specific individuals.
Consequently, explanation generation must provide meaningful
insight while satisfying strong privacy guarantees
\cite{TheAlgFoundDP2014}.

Existing work studied query refinement under diversity constraints
\cite{ERICA,MK09},
fairness-aware data processing and bias mitigation
\cite{MLJ22,LMSJ23,barocas-hardt-narayanan},
and database repairs and consistent query answering
\cite{arenas1999,bertossi2006,bertossi2011}.

The database community has also developed several approaches for
explaining query results and data transformations.
Examples include provenance-based explanations,
interactive explanation systems such as FEDEX \cite{FEDEX},
Tableau Explain Data \cite{TabEE},
and recent differentially private explanation frameworks such as
DPXPlain \cite{TaoGMR22}.

However, these approaches focus primarily on explaining a single
query result, a prediction, or a data transformation.
They do not address the problem of explaining the semantic
differences between two related query outputs, such as an original
query and its refined version or a query evaluated before and after
a database repair.

Furthermore, existing explanation systems typically assume direct
access to the underlying statistics or focus on releasing local
explanations.
To the best of our knowledge, no prior work provides a framework
for generating \emph{differentially private explanations of semantic
differences between related query outputs}.

In this paper, we address this gap by introducing a framework for
generating differentially private explanations of query-output
differences.
Rather than explaining a query result in isolation, we focus on
explaining the semantic changes between two related outputs,
such as an original query and its refined version or a query
evaluated before and after a database repair.

Our framework identifies informative explanations through a
two-stage process.

First, it privately selects promising semantic regions of the
data according to two complementary criteria:

\begin{itemize}

\item \textbf{Representativeness},
which favors regions supported by statistically meaningful
portions of the underlying dataset.

\item \textbf{Coverage},
which measures the amount of query-output evidence contained in
a region and favors regions supported by substantial numbers of
tuples from the compared query outputs.

\end{itemize}

The framework then privately selects explanations from the chosen
regions according to two additional criteria:

\begin{itemize}

\item \textbf{Interestingness},
which measures the extent to which an attribute distribution
differs between the two query outputs.

\item \textbf{Diversity},
which encourages the selected explanations to highlight different
attributes and therefore provide multiple perspectives on the
observed query transformation.

\end{itemize}

The framework first privately selects promising semantic regions
of the data, then privately chooses explanation candidates using
their interestingness and diversity contribution, and finally
releases noisy histogram summaries using standard differential
privacy mechanisms \cite{TheAlgFoundDP2014}.

The resulting explanations provide meaningful insight into the
effects of query refinements and repairs while preserving
individual privacy.

Although motivated by fairness-aware query refinement, the
proposed framework naturally generalizes to a broad range of
data-management tasks, including data cleaning, provenance
analysis, query debugging, and consistent query answering.
Our main contributions are as follows:

\begin{itemize}

\item We introduce the problem of generating differentially private
explanations for semantic differences between related query outputs.

\item We propose a view-based explanation model and a utility
framework that combines representativeness, coverage,
interestingness, and diversity.

\item We develop a differentially private explanation-generation
algorithm that combines private view selection, private explanation
selection, and noisy histogram release.

\item We provide formal privacy guarantees for the complete
explanation pipeline.

\item We experimentally evaluate the framework on multiple
real-world datasets and demonstrate that it produces informative
and interpretable explanations while preserving privacy.

\end{itemize}